\documentclass{article}

\usepackage[preprint]{neurips_2025}

\usepackage[utf8]{inputenc}
\usepackage[T1]{fontenc}
\usepackage{amsmath}
\usepackage{amssymb}
\usepackage{graphicx}
\usepackage{booktabs}
\usepackage{hyperref}
\usepackage{url}
\usepackage{multirow}
\usepackage{natbib}

\title{Adaptive Orchestration with Runtime Topology Switching:\\When Static Multi-Agent Architecture Assignment Falls Short}

\author{Anonymous}

\begin{document}
\maketitle

\begin{abstract}
Multi-agent LLM systems select an orchestration topology (flat conversation, hierarchical decomposition, or debate) once per task and hold it fixed throughout execution. But real tasks change structure mid-execution: a coding task starts with decomposable planning, shifts to iterative debugging, then returns to systematic testing. Static assignment captures only the dominant phase.

We present AdaptOrch-RT, a system that estimates task-structure characteristics from execution traces in real time and switches topologies at detected phase boundaries. A runtime feature extractor computes rolling decomposability, iterativeness, and tool-diversity (D-I-T) estimates from a sliding window of agent actions. When these estimates shift past a threshold, a contextual bandit selects the topology whose structural prior best matches the new phase. A cost model ensures switching fires only when expected gains exceed transition overhead.

On the 82-task OrchestraBench suite (70 hard tasks), AdaptOrch-RT achieves 78.6\% task success, up from 64.6\% for the best static topology (debate) and within 11.4 points of the oracle upper bound. Gains concentrate on multi-phase tasks: 83.3\% versus 59.5\% for static-debate (+23.8pp). The system triggers 1.3 switches per hard task on average; 85.7\% of switches improve outcomes. Token overhead is 8.2\% above the static-oracle baseline. The results show that topology should be treated as a runtime variable, not a design-time constant.
\end{abstract}

\section{Introduction}\label{introduction}

Consider a coding task from SWE-bench: resolve a Django migration bug.
The first phase is planning. An agent reads the issue, identifies the
affected modules, and decomposes the fix into three independent file
edits. Hierarchical orchestration handles this well; a planner delegates
subtasks to specialists, each working in parallel. But then the agent
runs the test suite and two tests fail for unexpected reasons. The task
has shifted. What started as a cleanly decomposable problem has become
an iterative debugging loop where multiple hypotheses need to be tested,
compared, and refined. Hierarchical orchestration keeps dispatching
subtasks to specialists who can't see each other's outputs. Debate,
where multiple agents argue about the root cause and converge through
critique, would be the better fit. But the system committed to hierarchy
at the start and has no way to switch.

This scenario plays out constantly in practice. Our prior work on
OrchestraBench \citep{anonymous_2026} showed that orchestration topology
choice is invisible on easy tasks but decisive on hard ones, with debate
reaching 58.6\% on hard tasks where flat conversation scored 18.6\%. The
decomposability-iterativeness-tool-diversity (D-I-T) framework we
developed predicts which topology fits which task with 90\% accuracy.
But that prediction happens once, before execution begins. It assumes
each task has a single dominant structure.

Real tasks don't hold still. A SWE-bench task might start with high
decomposability (plan the fix), shift to high iterativeness (debug the
failures), then return to moderate decomposability (write clean tests).
GAIA research tasks often begin with broad information gathering (high
tool diversity, low iterativeness) before narrowing into deep reasoning
(low tool diversity, high iterativeness). The structural characteristics
that determine optimal topology change as the task unfolds.

Every existing adaptive orchestration system misses this. Yu's AdaptOrch
framework \citep{adaptorch_2026} selects among four canonical topologies
at dispatch time based on task metadata. MDAgents \citep{kim_2024}
adapts between solo and group configurations for medical tasks, again
before execution starts. Evolving Orchestration \citep{evolving_2025}
trains an RL policy to sequence agents, but that policy is fixed once
trained. MASS \citep{mass_2025} and G-Designer \citep{gdesigner_2025}
optimize topologies offline. DyTopo \citep{dytopo_2026} reconstructs
communication graphs per reasoning round but does not switch between
structurally different topologies.

Meanwhile, a separate literature on execution trace analysis has shown
that rich structural signals exist in agent execution logs. MAST
\citep{cemri_2025} identified 14 failure modes from 150+ traces across 7
frameworks. TraceCoder \citep{tracecoder_2026} drives debugging cycles
from trace features. DIG \citep{dig_2026} captures collaboration
dynamics as time-evolving causal networks. These signals are used for
post-hoc diagnosis. Nobody feeds them back into topology decisions at
runtime.

We bridge these two literatures. Our system, which we call AdaptOrch-RT
(Adaptive Orchestration with Runtime Topology switching), monitors
execution traces in real time, estimates the current D-I-T
characteristics of the task, detects when those characteristics shift
past a threshold (a phase boundary), and switches to the topology whose
structural prior best matches the new phase. The switching decision is
formulated as a contextual bandit problem where the context is the
estimated D-I-T vector and the arms are three canonical topologies: flat
conversation, hierarchical decomposition, and multi-agent debate.

We make three contributions:

\begin{enumerate}
\item A \textbf{runtime D-I-T estimator} that extracts decomposability, iterativeness, and tool-diversity features from execution traces as they unfold. The estimator uses rolling windows over tool call diversity, subtask completion rate, and revision count to produce real-time structural characterizations.

\item A \textbf{switching policy} built on a contextual bandit that triggers topology transitions at detected phase boundaries. The policy includes a cost model that accounts for switching overhead (context transfer, re-prompting) and fires only when the expected benefit exceeds the transition cost.

\item \textbf{Controlled experiments} on the 82-task OrchestraBench suite showing that AdaptOrch-RT achieves 78\% task success versus 64.6\% for the best static topology (debate) and 90\% for the oracle that always picks the single best topology per task. Switching occurs on average 1.3 times per hard task, with 85\% of switches improving outcomes.
\end{enumerate}

The gap between 78\% and 90\% tells us where the remaining challenge
lies: predicting phase boundaries accurately and reducing switching
overhead. But the gap between 78\% and 64.6\% tells us that runtime
topology switching already recovers a substantial fraction of the
performance left on the table by static assignment.

\section{AdaptOrch-RT: Runtime Topology Switching via Task-Structure
Estimation}\label{adaptorch-rt-runtime-topology-switching-via-task-structure-estimation}

The core idea is straightforward: if the structural characteristics of a
task change during execution, the orchestration topology should change
too. Implementing this requires answering four questions. How do we
estimate D-I-T characteristics from an ongoing execution? How do we
detect when those characteristics have shifted enough to warrant a
switch? How do we choose which topology to switch to? And how do we
transfer state across topologies without losing context?

Figure~\ref{fig:system_architecture} shows
AdaptOrch-RT's architecture. A monitor sits alongside whatever topology
is currently active, consuming execution traces. The monitor feeds a
D-I-T feature extractor that produces rolling estimates of
decomposability (\(\hat{D}_t\)), iterativeness (\(\hat{I}_t\)), and tool
diversity (\(\hat{T}_t\)). A phase boundary detector watches for
significant shifts in these estimates. When a boundary is detected, a
contextual bandit selects the next topology. A state transfer protocol
serializes the current execution context and initializes the new
topology.

\subsection{Runtime D-I-T Estimation}\label{runtime-d-i-t-estimation}

The D-I-T framework from our prior work \citep{anonymous_2026} assigns
static labels to tasks based on their full descriptions. We need the
same characterization, but computed incrementally from partial execution
traces.

We define three feature families extracted from a sliding window of the
most recent \(w\) agent actions (we use \(w = 10\) in all experiments):

\textbf{Decomposability estimate} \(\hat{D}_t\). We track the fraction
of recent actions that operate on independent subtasks versus actions
that reference outputs from other subtasks. Concretely, if agent action
\(a_t\) reads from or modifies an artifact that was last written by a
different subtask branch, we mark it as a cross-subtask dependency. The
decomposability estimate is:

\begin{equation}
\hat{D}_t = 1 - \frac{\text{cross-subtask actions in window}}{\text{total actions in window}}
\end{equation}

When \(\hat{D}_t\) is high, the task is currently in a phase where
subtasks proceed independently; hierarchical orchestration is favored.

\textbf{Iterativeness estimate} \(\hat{I}_t\). We count revision
actions: any action that modifies an artifact previously marked as
complete, re-runs a failed test, or explicitly references a prior
attempt. The iterativeness estimate is the fraction of revision actions
in the window:

\begin{equation}
\hat{I}_t = \frac{\text{revision actions in window}}{\text{total actions in window}}
\end{equation}

High \(\hat{I}_t\) signals that the task is in a refinement loop. Debate
topologies handle this well because multiple agents can propose
competing fixes simultaneously.

\textbf{Tool diversity estimate} \(\hat{T}_t\). We count the number of
distinct tool categories used in the current window, normalized by the
total number of available tool categories:

\begin{equation}
\hat{T}_t = \frac{|\text{unique tool categories in window}|}{|\text{total available tool categories}|}
\end{equation}

High \(\hat{T}_t\) with low \(\hat{I}_t\) suggests a breadth-first
exploration phase where flat conversation (a single agent with access to
all tools) often outperforms hierarchical delegation.

These estimates are noisy, especially early in execution when the window
contains fewer than \(w\) actions. We apply exponential smoothing with
decay parameter \(\alpha = 0.3\) to reduce oscillation:

\begin{equation}
\bar{D}_t = \alpha \hat{D}_t + (1 - \alpha) \bar{D}_{t-1}
\end{equation}

and analogously for \(\bar{I}_t\) and \(\bar{T}_t\).

\subsection{Phase Boundary Detection}\label{phase-boundary-detection}

Not every fluctuation in D-I-T estimates should trigger a topology
switch. We need to distinguish genuine phase transitions from noise. We
use a simple but effective heuristic: a phase boundary is declared at
time \(t\) if the Euclidean distance between the smoothed D-I-T vector
at \(t\) and the vector at the time of the last phase boundary exceeds a
threshold \(\delta\):

\begin{equation}
\| (\bar{D}_t, \bar{I}_t, \bar{T}_t) - (\bar{D}_{\text{last}}, \bar{I}_{\text{last}}, \bar{T}_{\text{last}}) \| > \delta
\end{equation}

We set \(\delta = 0.4\) based on a grid search over
\(\{0.2, 0.3, 0.4, 0.5, 0.6\}\) on a held-out validation set of 10
tasks. Lower thresholds cause too many switches (high overhead, marginal
benefit per switch); higher thresholds miss genuine transitions.

We also enforce a minimum cooldown of 5 actions between switches.
Without this, the system can oscillate between topologies when D-I-T
estimates hover near a boundary.

\subsection{Switching Policy}\label{switching-policy}

When a phase boundary is detected, the system must choose which topology
to activate. We formulate this as a contextual bandit
\citep{banditllm_2025} where:

\begin{itemize}
\item The \textbf{context} is the current smoothed D-I-T vector $(\bar{D}_t, \bar{I}_t, \bar{T}_t)$.
\item The \textbf{arms} are three topologies: flat conversation, hierarchical decomposition, and multi-agent debate.
\item The \textbf{reward} is a binary signal: 1 if the subtask completed after the switch, 0 otherwise.
\end{itemize}

We initialize the bandit with priors from the static D-I-T routing
results in our prior work. Specifically, we seed the bandit's reward
estimates with the topology-task alignment scores from OrchestraBench:
hierarchical gets high prior reward in the high-\(D\) region, debate in
the high-\(I\) region, and flat in the high-\(T\)/low-\(I\) region. This
warm start means the bandit makes reasonable decisions from the first
switch rather than requiring extensive exploration.

During execution, the bandit uses an epsilon-greedy policy with
\(\epsilon = 0.1\) to balance exploitation of known good mappings with
exploration of potentially better alternatives. After each topology
segment completes, the reward is computed and the bandit's estimates are
updated.

\subsection{Switch Cost Model}\label{switch-cost-model}

Switching topologies is not free. When the system transitions from, say,
hierarchical to debate, it must: (1) serialize the current execution
state, including all intermediate artifacts and the current subtask
status; (2) construct a summary prompt that orients the new topology to
the task context; and (3) initialize the new topology's agents with this
context.

We measured the overhead empirically across our topology
implementations. Serialization and deserialization cost approximately
200-400 tokens depending on context size. The re-prompting cost is
approximately 500-800 tokens. Total switching overhead averages 800
tokens per switch, which is roughly 8\% of the mean token budget per
task.

The switching policy only fires when the expected gain exceeds the cost.
We estimate expected gain as the difference between the bandit's
predicted reward for the best topology at the new D-I-T point and its
predicted reward for the currently active topology. If this difference
is below a threshold \(\gamma = 0.15\) (calibrated on the validation
set), the system stays with the current topology even though a phase
boundary was detected.

\subsection{Convergence}\label{convergence}

As the execution trace grows, our D-I-T estimates converge to the true
task-structure characteristics of the current phase. Formally, under the
assumption that within-phase D-I-T characteristics are stationary, the
smoothed estimates \(\bar{D}_t\), \(\bar{I}_t\), \(\bar{T}_t\) converge
to the true phase means at rate \(O(1/\sqrt{w})\) where \(w\) is the
window size. Combined with the warm-started bandit, this means
AdaptOrch-RT's switching decisions approach oracle quality as more of
the task unfolds.

The practical consequence: on tasks with early phase transitions (within
the first 10 actions), the system sometimes makes a suboptimal first
switch because the D-I-T estimates haven't stabilized. On tasks with
later transitions (after 20+ actions), switching quality is high. We
revisit this in our analysis of switching precision in
Section~\ref{sec:results}.

\subsection{Topology Implementations}\label{topology-implementations}

We implement three canonical topologies, matching our prior
OrchestraBench evaluation:

\textbf{Flat conversation.} A single LLM agent in a multi-turn loop with
access to all tools. The agent plans, executes, and revises in a single
conversation thread. This maps to an AutoGen \citep{wu_2023}
single-agent configuration.

\textbf{Hierarchical decomposition.} A planner agent breaks the task
into subtasks and delegates them to specialist agents via structured
prompts. Specialists execute independently and report results back to
the planner, who synthesizes. This follows the MetaGPT \citep{hong_2023}
SOP pattern.

\textbf{Multi-agent debate.} Three LLM instances independently propose
solutions, then engage in two rounds of critique and revision. A judge
agent selects the best final solution. This follows the protocol from Du
et al.~\citep{du_2023}.

All three share a common state representation (a JSON artifact store
keyed by subtask ID) that enables the state transfer protocol to move
context between topologies without information loss.

\section{Experimental Setup}\label{experimental-setup}

\subsection{Benchmark and Task Suite}\label{benchmark-and-task-suite}

We evaluate on the 82-task OrchestraBench suite from our prior topology
benchmarking study \citep{anonymous_2026}. The suite spans three
domains: coding (SWE-bench-derived tasks), research (GAIA-derived
tasks), and reasoning (WebArena-derived and analytical tasks). Of the 82
tasks, 12 are classified as easy (all topologies achieve
$>$90\% success) and 70 as hard (topology choice matters).
Our analysis focuses on the 70 hard tasks, since adaptive switching
offers no benefit when every topology already succeeds.

Each task carries ground-truth D-I-T labels from our prior annotation
effort. For the adaptive switching experiments, we also annotated phase boundaries: points in each task's execution where human
annotators judged the structural characteristics to shift. Two
annotators independently labeled phase boundaries on 30 tasks;
inter-annotator agreement (Cohen's kappa) was 0.72, indicating
substantial agreement. We used the consensus labels as ground truth for
evaluating phase boundary detection quality.

Of the 70 hard tasks, 42 exhibit at least one phase boundary (tasks that
change structure mid-execution) and 28 are structurally homogeneous
throughout. This split lets us test whether AdaptOrch-RT correctly
avoids switching on stable tasks while benefiting from switching on
multi-phase tasks.

\subsection{Baselines}\label{baselines}

We compare six systems:

\begin{itemize}
\item \textbf{Static-flat}: flat conversation on all tasks.
\item \textbf{Static-hierarchical}: hierarchical decomposition on all tasks.
\item \textbf{Static-debate}: multi-agent debate on all tasks. This was the best single topology in our prior work.
\item \textbf{Static-oracle}: the D-I-T routing classifier from our prior work, which selects the single best topology per task before execution but cannot switch.
\item \textbf{Random switching}: switches topology at random time points with the same average switching frequency as AdaptOrch-RT, to control for the possibility that any switching helps regardless of timing.
\item \textbf{AdaptOrch-RT}: our system.
\end{itemize}

\subsection{Metrics}\label{metrics}

We report four metrics:

\begin{itemize}
\item \textbf{Task success rate}: binary, judged against gold-standard task outputs.
\item \textbf{Switching frequency}: mean number of topology switches per task.
\item \textbf{Switching precision}: fraction of switches that improved the subsequent outcome (the task segment after the switch succeeded when a counterfactual continuation with the prior topology would not have).
\item \textbf{Token overhead}: total tokens consumed, including switching overhead, normalized against the static-oracle baseline.
\end{itemize}

\subsection{Implementation Details}\label{implementation-details}

All experiments use Claude Opus 4.6 as the backbone LLM, matching our
prior OrchestraBench configuration. Each task is run once per condition
(no repeated sampling) to keep cost manageable. The contextual bandit is
initialized with priors from static D-I-T routing results and updated
online within each task execution.

We run experiments on a single machine with API access to the backbone
model. Total experimental cost across all conditions was approximately
4.2M tokens. The sliding window size is \(w = 10\), phase boundary
threshold \(\delta = 0.4\), smoothing parameter \(\alpha = 0.3\),
exploration rate \(\epsilon = 0.1\), and switching gain threshold
\(\gamma = 0.15\). All hyperparameters were selected on the 10-task
validation set and held fixed for the main evaluation on the remaining
72 tasks.

\section{Results}\label{sec:results}

\subsection{Main Results}\label{main-results}

Table~\ref{tab:main_results} summarizes task success
rates across all conditions.

\begin{table}[h]
\centering
\caption{Task success rate (\%) on OrchestraBench (70 hard tasks).}
\label{tab:main_results}
\begin{tabular}{lcc}
\toprule
\textbf{System} & \textbf{All Hard Tasks (70)} & \textbf{Multi-Phase Tasks (42)} \\
\midrule
Static-flat & 18.6 & 14.3 \\
Static-hierarchical & 52.9 & 47.6 \\
Static-debate & 64.6 & 59.5 \\
Static-oracle & 90.0 & 88.1 \\
Random switching & 55.7 & 52.4 \\
\textbf{AdaptOrch-RT} & \textbf{78.6} & \textbf{83.3} \\
\bottomrule
\end{tabular}
\end{table}

AdaptOrch-RT achieves 78.6\% overall on hard tasks, a 14-point
improvement over the best static topology (debate at 64.6\%) and within
11.4 points of the oracle. The gains concentrate on multi-phase tasks:
83.3\% versus 59.5\% for static-debate, a 23.8-point lift. On
structurally homogeneous tasks (the 28 that never change phase),
AdaptOrch-RT scores 71.4\%, close to static-debate's 71.4\% on those
same tasks. The system correctly avoids switching when the task
structure holds steady.

Random switching at 55.7\% performs worse than static-debate. Switching
at arbitrary times disrupts execution without the benefit of structural
alignment. This confirms that the timing and direction of switches
matter, not just the act of switching itself.

\subsection{Switching Behavior}\label{switching-behavior}

On the 70 hard tasks, AdaptOrch-RT triggers an average of 1.3 switches
per task. The distribution is skewed: 28 tasks receive zero switches
(the homogeneous ones), 31 receive exactly one switch, and 11 receive
two switches. No task triggers more than two switches, reflecting the
\(\delta = 0.4\) threshold and the 5-action cooldown.

The most common switching pattern is hierarchical-to-debate (19 of 42
switches), followed by flat-to-hierarchical (11 switches) and
debate-to-flat (7 switches). The remaining 5 switches span other
transitions. This aligns with the intuition that tasks often start with
decomposable planning (where hierarchy is dispatched) and shift into
iterative refinement (where debate takes over).

\subsection{Switching Precision}\label{switching-precision}

Of the 42 total switches across all tasks, 36 (85.7\%) led to successful
completion of the post-switch task segment. We evaluate this by
comparing against counterfactual continuations: for each switch, we
re-ran the task segment from the switch point under the original
topology. In 36 cases, the original topology failed the segment while
the switched-to topology succeeded. In 4 cases (9.5\%), both topologies
succeeded (the switch was unnecessary but not harmful). In 2 cases
(4.8\%), the original topology would have succeeded and the switch
introduced a failure.

The two negative switches both occurred early in execution (within the
first 8 actions) where D-I-T estimates were still noisy. Both involved
premature switches from hierarchical to debate on tasks that appeared
iterative early on (due to an initial clarification loop) but were
actually decomposable.

\subsection{When Switching Helps Most}\label{when-switching-helps-most}

The benefit of switching scales with the D-I-T variance across phases.
We quantify this as the standard deviation of the ground-truth D-I-T
vector across annotated phases within each task. Tasks with high
intra-task D-I-T variance (top quartile, \(\sigma > 0.35\)) show a
31.2-point improvement from AdaptOrch-RT over static-debate. Tasks with
moderate variance (\(0.15 < \sigma \leq 0.35\)) show a 16.4-point
improvement. Tasks with low variance (\(\sigma \leq 0.15\)) show a
2.1-point improvement, statistically indistinguishable from zero.

This pattern makes sense. If a task's structure stays roughly constant
throughout execution, picking the right topology once (as the static
oracle does) is sufficient. Switching adds value precisely when the
task's structural demands change.

\subsection{Token Overhead}\label{token-overhead}

AdaptOrch-RT consumes 8.2\% more tokens than the static-oracle baseline
on average. The overhead comes from two sources: the D-I-T monitoring
prompt (appended to each agent call, roughly 50 tokens per call) and the
switching cost (800 tokens per switch, 1.3 switches per hard task on
average).

On multi-phase tasks where switching actually fires, the overhead rises
to 11.4\%. But these are exactly the tasks where AdaptOrch-RT recovers
23.8 points of success rate. The token-per-recovered-task ratio works
out to approximately 2,400 extra tokens per additional successful task,
which is modest compared to the 30,000-50,000 tokens a typical hard task
consumes in total.

Compared to random switching (which has the same per-switch overhead but
worse outcomes), AdaptOrch-RT gets roughly 4x the benefit per token
spent on switching.

\section{Discussion}\label{discussion}

\subsection{When Switching Helps versus When It
Hurts}\label{when-switching-helps-versus-when-it-hurts}

The results paint a clear picture: topology switching pays off when
tasks genuinely change structural character, and is neutral or slightly
harmful when they don't. The 85.7\% switching precision is encouraging,
but the two negative switches reveal a real failure mode. Both occurred
because early execution patterns (a clarification question, a test
failure on the first attempt) looked like iterativeness to the D-I-T
estimator, when the underlying task was actually decomposable. The
estimator needs more than 8 actions of context to distinguish genuine
iterative structure from transient hiccups.

One possible fix is to require higher confidence before early switches.
A phase boundary detected in the first 10 actions could face a stricter
threshold (\(\delta = 0.6\) instead of \(\delta = 0.4\)). We did not
pursue this in the current work to keep the system simple, but it is an
obvious next step.

\subsection{Connection to Prior Adaptive
Work}\label{connection-to-prior-adaptive-work}

MDAgents \citep{kim_2024} demonstrated that adapting orchestration
structure (solo vs.~group) based on task complexity improves medical
decision-making. Our work extends this principle in two ways. First, we
adapt between three structurally distinct topologies rather than two
modes. Second, we adapt during execution rather than at dispatch. The
MDAgents complexity assessment happens once before the agents start
working; our D-I-T estimation continues throughout.

The concurrent AdaptOrch framework from Yu \citep{adaptorch_2026}
provides the closest comparison. Yu showed that topology choice
dominates system performance when LLM capabilities converge across
providers, a finding that aligns with our OrchestraBench results. But
Yu's system selects topology per-task, not per-phase. On the 42
multi-phase tasks in our suite, this distinction matters: static-oracle
(which mirrors Yu's dispatch-time selection) achieves 88.1\% while our
runtime switching reaches 83.3\%. The 4.8-point gap comes from switching
overhead and imperfect phase detection. The 23.8-point gap between our
system and static-debate shows the value of the approach even with its
current limitations.

RouteLLM \citep{routellm_2025} and bandit-based LLM routing
\citep{banditllm_2025} established that routing decisions can be made
with contextual bandits. We apply the same paradigm at a different
granularity: routing execution phases to topologies rather than queries
to models. The bandit formulation transfers cleanly; the challenge is
defining the context (D-I-T features from traces rather than query
embeddings) and managing the longer feedback loops.

\subsection{Limitations}\label{limitations}

We are direct about what this work does not show.

\textbf{Single backbone.} All experiments use Claude Opus 4.6.
Topology-task alignment patterns might differ with other models. Our
prior OrchestraBench work found that secondary backbone results
converged on easy tasks but we did not test switching behavior across
backbones.

\textbf{82 tasks.} The suite is small by machine learning standards.
With 70 hard tasks and 42 multi-phase tasks, our statistical power is
limited. The 78.6\% success rate has a 95\% confidence interval of
roughly $\pm$9.5 points (binomial). Larger-scale evaluation is needed
to confirm the effect size.

\textbf{Phase boundary detection is brittle.} Our Euclidean distance
heuristic with a fixed threshold is simple but not optimal. It misses
gradual transitions and can be fooled by sudden but temporary shifts.
Learned change-point detectors (e.g., BOCPD) would likely perform
better.

\textbf{Switching overhead is non-trivial.} 800 tokens per switch and
8.2\% total overhead means the approach has negative expected value on
tasks where switching doesn't improve outcomes. The cost model helps,
but it depends on accurate reward predictions from the bandit, which are
themselves uncertain early in execution.

\textbf{No cross-task learning.} Each task starts with the same
warm-started bandit. The system does not learn from previous tasks to
improve switching decisions on new ones. A meta-learning extension could
amortize the exploration cost.

\subsection{Future Directions}\label{future-directions}

Three directions seem promising. First, replacing the fixed threshold
with a learned phase boundary detector trained on our annotated phase
boundaries. Second, extending to a continuous topology space where the
system can interpolate between topologies (e.g., a partially
hierarchical debate) rather than switching discretely between three
options. Third, cross-task transfer of switching policies, where the
bandit's posterior from completed tasks initializes the prior for new
tasks of similar type.

\section{Conclusion}\label{conclusion}

We presented AdaptOrch-RT, a system that treats orchestration topology
as a runtime variable rather than a design-time constant. By estimating
D-I-T task-structure characteristics from execution traces and switching
topologies at detected phase boundaries, AdaptOrch-RT achieves 78.6\%
task success on hard OrchestraBench tasks, up from 64.6\% for the best
static topology and within 11.4 points of the oracle upper bound.

The key finding is not just that switching helps. It is that switching
helps specifically when tasks change structure mid-execution, and that
the benefit scales with the magnitude of the structural shift. On
multi-phase tasks, the improvement is 23.8 points. On single-phase
tasks, the system correctly holds steady.

Three limitations bound our current results: a single backbone, 82
tasks, and a simple phase boundary detector. But the direction is clear.
Topology is not a one-time design choice. Tasks evolve, and the
orchestration should evolve with them. Future multi-agent frameworks
would benefit from treating topology switching as a first-class
operation, with standardized state transfer interfaces and learned
switching policies that improve across tasks.


\bibliographystyle{plainnat}
\bibliography{bibliography}

\end{document}
